\section{Теоретические основы визуальной навигации}

В части о зрении вы уже находили метки ArUco, выбирали нужный ID, считали центр метки и ошибку относительно центра кадра, садились на метку. Эта часть объясняет, \emph{как} устроены эти шаги внутри детектора и что нужно добавить, чтобы перейти от пикселей к метрам.

\begin{infobox}[Где выполнять задания этой части]
\begin{itemize}
    \item \textbf{Симулятор.} Камера \verb|Camera(view="bottom")|, детектор \verb|cv2.aruco| со словарями OpenCV, номера и углы меток, наведение \verb|set_manual_speed|, телеметрия, процедурные полигоны с метками. Распознавание здесь моделируется по геометрии сцены: детектор возвращает углы тех меток, которые целиком видны камере. Внутренние шаги OpenCV (порог, контуры, \verb|solvePnP|, калибровка) в браузере не выполняются: \verb|cv2.solvePnP| в симуляторе всегда отвечает \verb|(False, None, None)|, а \verb|cv2.imwrite| ничего не сохраняет.
    \item \textbf{Компьютер с OpenCV} (\verb|pip install opencv-contrib-python|). Генерация маркеров, разбор конвейера детектора, фильтры, файлы карт, графики. Дрон для этого не нужен.
    \item \textbf{Настоящий дрон или веб-камера.} Калибровка камеры, \verb|solvePnP| и расстояние в метрах, печать и монтаж маркеров, освещение, блики.
\end{itemize}
Каждое задание ниже помечено, где его выполнять.
\end{infobox}

\begin{missionbox}
Перед этой главой пройдите в панели \ui{Обучение} миссии Python №\,73--81 (камера) и №\,82--90 (ArUco шаг за шагом). Ими проверяются словарь, детекция, отсутствие метки, четыре угла, центр и устойчивое распознавание --- всё, на что опирается эта глава.
\end{missionbox}

\subsection{Введение в ArUco и AprilTag}
Визуальная навигация — это способ определения положения робота в пространстве с помощью анализа изображений. В отличие от GPS, который работает только на улице и имеет погрешность в несколько метров, визуальные маркеры позволяют дрону позиционироваться с точностью до сантиметров даже в закрытых помещениях.

\textbf{ArUco} (Augmented Reality University of Cordoba)~\cite{garrido2014aruco} и \textbf{AprilTag}~\cite{olson2011apriltag} — это самые популярные семейства фидуциальных (опорных) маркеров. Они представляют собой черно-белые квадраты, которые содержат закодированную информацию (ID маркера) и позволяют алгоритмам компьютерного зрения однозначно определить ориентацию маркера в пространстве. В OpenCV оба семейства распознаёт один и тот же модуль \verb|cv2.aruco|~\cite{opencv_aruco}: AprilTag в нём --- просто ещё один словарь.

\begin{taskbox}[Свой и чужой словарь]
\emph{Симулятор.} Над меткой 12 (сцена миссии Python №\,83) создайте два детектора: со словарём \verb|DICT_4X4_50| и со словарём \verb|DICT_6X6_250|. Обработайте ими один и тот же кадр и напечатайте оба списка ID.
\expected{«чужой» словарь $4\times4$ печатает пустой список, словарь \texttt{DICT\_6X6\_250} --- \texttt{[12]}.}
Учтите отличие модели: симулятор сравнивает словари по имени, поэтому и детектор \verb|DICT_6X6_50| метку из \verb|DICT_6X6_250| не найдёт. Настоящий OpenCV её найдёт: первые 50 маркеров этих словарей совпадают.
\mission{Python}{82--83}
\end{taskbox}

\subsection{Принцип работы алгоритма обнаружения}
Процесс обнаружения маркера на изображении состоит из нескольких этапов. Понимание этого конвейера (pipeline) важно для диагностики проблем, когда дрон <<не видит>> маркеры.

\subsubsection{1. Предварительная обработка (Preprocessing)}
Камера передает цветное изображение (в OpenCV — в порядке каналов BGR). Алгоритм сначала переводит его в оттенки серого (Grayscale), так как цвет не несет информации для черно-белых маркеров.

Затем применяется \textbf{адаптивная бинаризация} (Adaptive Thresholding). В отличие от простого порога (где все пиксели ярче 127 становятся белыми, а темнее — черными), адаптивный метод вычисляет порог для каждой небольшой области изображения. Это позволяет находить маркеры даже при неравномерном освещении (например, когда одна часть маркера в тени).

Программа ниже показывает весь конвейер на одной картинке. Без аргументов она сама рисует маркер ID~12 на листе, правая часть которого в тени, и сравнивает общий и адаптивный порог; с именем файла --- обрабатывает вашу фотографию. Запускайте её на компьютере: в симуляторе этих функций OpenCV нет.

\begin{codefiletitle}[python]{Конвейер детектора по шагам (OpenCV, компьютер)}{python/advanced_navigation/binarization_example.py}
\end{codefiletitle}

\subsubsection{2. Поиск контуров (Contour Extraction)}
Алгоритм ищет границы связных областей на бинарном изображении (после \verb|THRESH_BINARY_INV| чёрная рамка маркера становится белой областью). Из всех найденных контуров отбираются только те, которые:
\begin{itemize}
    \item Являются замкнутыми.
    \item Аппроксимируются четырехугольником (имеют 4 вершины).
    \item Являются выпуклыми (Convex).
    \item Имеют достаточный размер (площадь).
\end{itemize}
Именно это делает функция \verb|count_quads| в примере выше.

\subsubsection{3. Исправление перспективы (Perspective Warping)}
Найденный четырехугольник может быть искажен перспективой (выглядеть как трапеция или ромб), если камера смотрит на него под углом. Алгоритм вычисляет матрицу гомографии, чтобы <<выпрямить>> изображение и получить идеальный квадрат (Canonical Marker Image).

\subsubsection{4. Декодирование (Decoding)}
Выпрямленный квадрат разбивается на сетку: биты данных плюс чёрная рамка шириной в одну клетку (у маркера $6 \times 6$ это $8 \times 8$ клеток). Каждая клетка анализируется: черная — 0, белая — 1. Полученная матрица битов сверяется со словарем (Dictionary) во всех четырёх поворотах --- так детектор узнаёт и номер маркера, и то, какой из углов «первый». Поэтому \verb|corners| всегда перечисляет углы в одном порядке: по часовой стрелке, начиная с верхнего левого угла самого маркера.

\begin{taskbox}[Тень и порог]
\emph{Компьютер с OpenCV.} Запустите \verb|binarization_example.py| без аргументов. Сравните окна \texttt{global threshold} и \texttt{adaptive threshold} и числа четырёхугольников. Затем сделайте тень глубже: в \verb|make_test_image| замените \verb|0.25| на \verb|0.1|, --- и снова сравните.
\expected{при общем пороге белые клетки маркера в тени становятся «чёрными» и четырёхугольников меньше, адаптивный порог сохраняет рамку и клетки; детектор печатает \texttt{Найдены ID: [12]}. Объясните, почему от глубины тени зависит только общий порог.}
\end{taskbox}

\subsection{Коррекция ошибок и Расстояние Хэмминга}
Маркеры ArUco спроектированы так, чтобы быть устойчивыми к ошибкам считывания. Даже если несколько пикселей (битов) будут распознаны неверно (из-за блика или грязи), алгоритм все равно сможет узнать маркер.

Это достигается за счет избыточности информации. Каждый словарь имеет параметр \textbf{Hamming Distance} (Расстояние Хэмминга) — минимальное количество битов, на которое отличаются любые два маркера в словаре (с учётом их поворотов).
\begin{itemize}
    \item Если словарь имеет расстояние Хэмминга $d = 4$, система может обнаружить до $d-1 = 3$ ошибок и исправить до $\lfloor (d-1)/2 \rfloor = 1$ ошибки.
    \item Чем больше бит в маркере (например, $6 \times 6$ против $4 \times 4$), тем дальше друг от друга можно разнести маркеры словаря и тем ниже риск ложных срабатываний. Чем больше маркеров в словаре при той же сетке, тем меньше это расстояние: \verb|DICT_6X6_50| надёжнее, чем \verb|DICT_6X6_1000|.
\end{itemize}

\begin{taskbox}[Расстояние Хэмминга своими руками]
\emph{Компьютер или симулятор, только Python.} Запишите два маркера $4 \times 4$ списками строк: \verb|A = ["1010", "0110", "1100", "0011"]| и \verb|B = ["1011", "0100", "1100", "1011"]|. Напишите функцию \verb|hamming(a, b)| (число несовпадающих битов) и функцию \verb|rotate(m)| (поворот матрицы на 90° по часовой стрелке). Напечатайте расстояние от \verb|A| до \verb|B| и до трёх поворотов \verb|B|.
\expected{без поворота --- 3; для поворотов на 90°, 180° и 270° --- 7, 11 и 9. Минимум (3) и есть расстояние между этими маркерами с учётом поворотов: такой словарь обнаружит 2 ошибки и исправит одну.}
\end{taskbox}

\subsection{Сравнение типов словарей (Dictionaries)}
При выборе словаря нужно искать баланс между количеством доступных ID и размером маркера (читаемостью).

\begin{table}[H]
\centering
\begin{tabularx}{\textwidth}{|l|X|X|X|}
\hline
\textbf{Словарь} & \textbf{Размер сетки} & \textbf{Кол-во ID} & \textbf{Применение} \\ \hline
DICT\_4X4\_50 & $4 \times 4$ бита & 50 & Дальние дистанции, низкое разрешение. Меньше деталей — легче распознать. \\ \hline
DICT\_5X5\_1000 & $5 \times 5$ бит & 1000 & Компромисс: много номеров при средней надёжности. \\ \hline
DICT\_6X6\_250 & $6 \times 6$ бит & 250 & Высокая надежность, используется в картах Пионера и в симуляторе; его первые 50 маркеров --- это \texttt{DICT\_6X6\_50}. \\ \hline
DICT\_APRILTAG\_36h11 & $6 \times 6$ бит & 587 & Промышленный стандарт, очень высокая устойчивость к ложным срабатываниям (минимальное расстояние Хэмминга 11). \\ \hline
\end{tabularx}
\caption{Характеристики словарей маркеров}
\end{table}

\subsection{Математика позиционирования (PnP Problem)}
После того как маркер найден и идентифицирован, нужно определить положение камеры относительно него. Эта задача называется \textbf{Perspective-n-Point (PnP)}.

Мы знаем:
\begin{enumerate}
    \item Координаты 4 углов маркера в его собственной системе координат (3D object points), например $(-S/2, S/2, 0)$, $(S/2, S/2, 0)$, $(S/2, -S/2, 0)$, $(-S/2, -S/2, 0)$, где $S$ — размер маркера, а начало координат --- в его центре. Порядок точек должен совпадать с порядком углов, который возвращает детектор.
    \item Координаты этих же 4 углов на изображении камеры (2D image points) в пикселях.
    \item Внутренние параметры камеры (Camera Matrix и Distortion Coefficients), полученные при калибровке~\cite{zhang2000calibration}.
\end{enumerate}

Функция \texttt{cv2.solvePnP} решает систему уравнений и возвращает признак успеха и два вектора:
\begin{itemize}
    \item \texttt{rvec} (Rotation Vector) — вектор вращения (в формате Родригеса), показывающий наклон камеры.
    \item \texttt{tvec} (Translation Vector) — вектор смещения $(x, y, z)$, показывающий позицию маркера относительно центра камеры.
\end{itemize}

\begin{alertbox}[Важно!]
Вектор \texttt{tvec} показывает координаты \textit{маркера в системе координат камеры}. Чтобы получить координаты \textit{дрона в системе координат маркера} (для навигации), нужно инвертировать эту трансформацию:
$$ R = Rodrigues(rvec) $$
$$ Position_{drone} = -R^T \cdot tvec $$
\end{alertbox}

\begin{warnbox}[Что работает в симуляторе]
В симуляторе \verb|cv2.solvePnP| --- заглушка: она возвращает \verb|(False, None, None)|, потому что у модели нет настоящей оптики и калибровки. Метры в симуляторе берут из телеметрии (\verb|get_local_position_lps|), а метку используют как пиксельную цель: ошибку центра метки относительно центра кадра (миссии 89 и 94). Расчёт \verb|tvec| проверяйте на настоящей камере.
\end{warnbox}

\begin{taskbox}[Масштаб вместо solvePnP]
\emph{Симулятор.} Модель камеры-«обскуры» говорит: видимая сторона метки в пикселях $s \approx f \cdot S / d$, где $d$ --- расстояние до метки. Значит, произведение $s \cdot d$ не должно зависеть от высоты. Зависните над меткой 12 на высотах 2, 3 и 4~м (у метки на полу $d$ равна высоте полёта $z$). На каждой высоте по свежему кадру посчитайте среднюю длину четырёх сторон метки по её углам (\verb|math.dist| между соседними углами) и напечатайте $z$, $s$ и $s \cdot z$.
\expected{три строки, в которых $s$ уменьшается с высотой, а $s \cdot z$ почти одинаково (расхождение в пределах нескольких процентов). Отсюда, зная $s \cdot z$ для одной высоты, можно оценить высоту по размеру метки.}
\mission{Python}{78, 87--88}
\end{taskbox}

\begin{taskbox}[Расстояние по tvec]
\emph{Настоящий дрон или веб-камера.} Откалибруйте камеру (глава «Создание навигационных карт»), распечатайте маркер со стороной 10~см и запустите пример определения координат маркера из главы об ArUco. Поставьте маркер на 0,5, 1,0 и 1,5~м от камеры (рулеткой) и запишите третью компоненту \verb|tvec|.
\expected{\texttt{tvec[2]} отличается от измеренного рулеткой расстояния не больше чем на 5\%. Если ошибка растёт с расстоянием пропорционально --- в коде указан неверный размер маркера; если она хаотична --- переделайте калибровку.}
\end{taskbox}

\subsection{Практические задания}

\begin{taskbox}[Анализ словарей]
\emph{Компьютер и принтер, камера дрона или веб-камера.} Сгенерируйте маркер ID=1 из словарей $4 \times 4$, $5 \times 5$ и $6 \times 6$ одного физического размера (например, 5~см) скриптом \verb|19_map_creation_example2_generation.py|.
\begin{enumerate}
    \item Распечатайте их.
    \item Отходите с камерой, пока маркеры распознаются.
    \item Какой маркер перестанет распознаваться первым? Почему?
\end{enumerate}
\expected{таблица «словарь --- предельное расстояние»; первым пропадает $6 \times 6$: при том же размере его клетки мельче и раньше сливаются в пиксели камеры.}
\end{taskbox}

\begin{taskbox}[Влияние освещения]
\emph{Настоящая камера.} Запустите \verb|binarization_example.py| с фотографией маркера или детектор из главы об ArUco на живом видео.
\begin{enumerate}
    \item Постепенно затемняйте комнату. На каком уровне освещенности маркер пропадет?
    \item Посветите фонариком прямо на маркер (создайте блик). Как это влияет на бинаризацию (посмотрите окно с ч/б изображением)?
\end{enumerate}
\expected{описание двух порогов: при темноте маркер пропадает из-за шума и смаза, при блике --- из-за белого пятна на месте чёрных клеток; в окне бинаризации видно, какие клетки «перевернулись».}
\end{taskbox}

\begin{questionbox}[Контрольные вопросы]
\begin{enumerate}
    \item Почему для навигации не используются QR-коды? (Подсказка: сравните сложность детектирования угла наклона).
    \item Что такое <<адаптивная бинаризация>> и зачем она нужна?
    \item Если дрон видит маркер под углом 45 градусов, как алгоритм понимает, что это квадрат, а не ромб?
    \item В чем физический смысл вектора \texttt{tvec}?
    \item Почему в симуляторе метку используют как пиксельную цель, а не как источник координат в метрах?
\end{enumerate}
\end{questionbox}
